\section{H3: Is the settlement price displaced?}\label{sec:distortion}

\subsection{Reversal of the stock settlement price}

Table~\ref{tab:reversal} estimates equation~\eqref{eq:reversal}. On ordinary sessions the move
over the window is not reversed the next morning: the coefficient for the securities with
derivatives is $\RevDerivativeDrift{}$, with a standard error of \RevDerivativeDriftSe{}. On
monthly expiries a part of the move is reversed. The additional coefficient is
$\RevDerivativeDriftxmonthlyexpiry{}$ (standard error \RevDerivativeDriftxmonthlyexpirySe{}),
so that about a quarter of the window's move is undone by the next morning, and the
randomization $p$-value is \RevDerivativeDriftxmonthlyexpiryRip{}. The corresponding
coefficients for weekly index expiries and month ends are smaller and not significant.

\begin{table}[htbp]
\centering
\small
\caption{Reversal of the settlement price. Regressions of the return from the settlement price
to the volume-weighted average price of the next session's first thirty minutes on the move
over the window and its interactions with the session types, with security and session fixed
effects. The last column is the difference between the securities with derivatives and the
placebo group. Standard errors clustered by session in parentheses. The randomization
$p$-value compares the monthly-expiry coefficient with its distribution over \RiDraws{}
samples in which the expiries are replaced by one ordinary session from each month.}
\label{tab:reversal}
\resizebox{\textwidth}{!}{\input{generated/tables/reversal}}
\end{table}

The reversal is not, however, confined to securities with an expiring contract. In the placebo
group the additional coefficient on monthly expiries is $\RevPlaceboDriftxmonthlyexpiry{}$
(randomization $p$-value \RevPlaceboDriftxmonthlyexpiryRip{}), and the difference between the
two sets of securities, $\RevDidDriftxmonthlyexpiry{}$, is close to zero (randomization
$p$-value \RevDidDriftxmonthlyexpiryRip{}). A settlement price displaced by trading aimed at
the settlement should be reversed where there is a settlement to influence and not elsewhere.
The pattern found here is a property of the monthly expiry session as a whole, shared by
securities on which nothing settles.

Figure~\ref{fig:path} shows the same result as a path. For each security and session the
prices through the window and at the next morning's open are measured relative to the average
of all securities on that session, and oriented in the direction of the window's move. On
ordinary sessions the move is kept overnight; on monthly expiries part of it is given back, to
a similar extent in all three groups.

\begin{figure}[htbp]
  \centering
  \includegraphics[width=\textwidth]{window_path.pdf}
  \caption{Average path of the price through the settlement window and at the next morning's
    open, relative to the pre-window price and to the average of all securities on the same
    session, oriented so that the window's move is positive.}
  \label{fig:path}
\end{figure}

On the order book sample the settlement price ends further from the last traded price on
expiry sessions, \AbsReversalBpsExpiry{} basis points against \AbsReversalBpsControl{}, and
less volume is traded within a basis point of the running average. Both are consistent with a
more volatile window rather than with trading that tracks the average. The first is larger in
securities with derivatives than in the placebo group, by \DidAbsreversalbpsPlacebo{} basis
points ($p = \DidAbsreversalbpsPlaceboP{}$). On the full-year sample, which covers every
session and uses session fixed effects, the corresponding comparison of the final minutes with
the settlement price is not specific to these securities (Section~\ref{sec:window}).

\subsection{The index, where settlement is in cash}\label{sec:index}

Settlement by delivery removes the incentive to move a stock's settlement price
(Section~\ref{sec:derivatives}). It does not remove it for the index. Nifty 50 futures and
options are settled in cash against the index's closing value, which is computed from the
closing prices of the constituents, that is from their settlement-window VWAPs. The index
options expired on every Thursday of the year and were the most heavily traded derivatives on
the exchange. If the settlement window is used to move a price, the index is where it pays.

The exchange's index file records, every second, the value of the Nifty 50 and of the Nifty
Next 50, an index of the next fifty large companies on which no derivatives were traded. The
second index is exposed to the same market and the same calendar but to no settlement, and it
serves as the placebo at the index level. The closing value of the Nifty 50 is approximated by
its average over the window; on the dates on which the published closing value was checked,
the approximation was within a few index points of it. The comparison uses all
\IdxSessions{} sessions, of which \IdxExpiries{} were index expiries.

\begin{table}[htbp]
\centering
\small
\caption{The index settlement. Regressions of each measure on indicators for index expiries,
every Thursday on which the index options expire, and for the monthly expiry in addition,
with a control for month ends. Difference is the Nifty 50 less the Nifty Next 50, which has no
derivatives. Heteroskedasticity-robust standard errors in parentheses. The randomization
$p$-value compares the index-expiry coefficient with its distribution over \RiDraws{} samples
in which each expiry is replaced by another session of the same week.}
\label{tab:index}
\resizebox{0.95\textwidth}{!}{\input{generated/tables/index}}
\end{table}

Table~\ref{tab:index} reports the results. Neither index, taken on its own, moves in a
consistent direction over the window on expiry days; the level of each is dominated by the
movement of the market as a whole. Their difference is not. On index expiries the change of the
Nifty 50 over the window, relative to the Nifty Next 50, differs from other sessions by $\IdxSpreadDriftExpiry{}$ basis points (standard error \IdxSpreadDriftExpirySe{}, randomization $p$-value
\IdxSpreadDriftExpiryRip{}). This is about \IdxSpreadDriftPoints{} index points at the year's
average level of \IdxNiftyLevel{}. It occurs on the weekly expiries, on which only the index
options settle, as much as on the monthly ones: the additional coefficient for monthly expiries
is $\IdxSpreadDriftMonthly{}$ and insignificant. The Nifty 50 ends the window below the Nifty
Next 50 on \IdxNegshareExpiry{}\% of index expiries, against \IdxNegshareOther{}\% of other
sessions. The movement takes place within the window: from its first minute to its last the
difference changes by $\IdxSpreadMoveExpiry{}$ basis points relative to other sessions
(randomization $p$-value \IdxSpreadMoveExpiryRip{}). Figure~\ref{fig:indexpath} shows the path.

\begin{figure}[htbp]
  \centering
  \includegraphics[width=0.72\textwidth]{index_path.pdf}
  \caption{The Nifty 50 relative to the Nifty Next 50 through the settlement window and at the
    next morning's open, relative to the pre-window level, on index expiries and on other
    sessions. Means with 95\% intervals.}
  \label{fig:indexpath}
\end{figure}

The next morning the difference recovers by \IdxSpreadNextrevExpiry{} basis points more than
after other sessions, which is of the same size as the fall and in the opposite direction. The
estimate is imprecise (randomization $p$-value \IdxSpreadNextrevExpiryRip{}), because overnight
returns vary far more than returns within the window, so the reversal is consistent with the
fall being temporary but does not establish it.

The same pattern appears in the stocks. On index expiries the liquid group, most of whose
members are constituents of the Nifty 50, changes from the first minute of the window to the last
by $\StockidxLiquidMoveOneFiveThreeZeroIndexexpiry{}$ basis points relative to the placebo group
(standard error \StockidxLiquidMoveOneFiveThreeZeroIndexexpirySe{}). The illiquid group, most of
whose members are outside the index, moves by $\StockidxIlliquidMoveOneFiveThreeZeroIndexexpiry{}$
basis points, which is not significant.

The index settlement is therefore displaced downward on expiry days, by an amount that is small
relative to the daily movement of the index but systematic across the year's expiries, specific
to the index that carries derivatives, and concentrated in the window in which its settlement
value is fixed. The data identify participant categories but not individual traders, and they
do not show the positions held in index options, so the cause cannot be established. Two
explanations are consistent with the evidence: trading intended to lower the settlement value,
and the unwinding at the settlement price of cash hedges held against index option positions.
Either way, the settlement value of the most traded derivative in the market is systematically
lowered in the window relative to an index without derivatives.
